\documentclass[11pt]{article}
\usepackage[a4paper,margin=1in]{geometry}
\usepackage{amsmath,amssymb}
\usepackage{booktabs}
\usepackage{siunitx}
\usepackage{hyperref}

\title{Independent Replication Report --- OSTI 3023480}
\author{X-100 Replication Project (subagent OSTI-3023480, wave 2026-07-01)}
\date{2026-07-05}

\begin{document}
\maketitle

\begin{abstract}
Huang \emph{et al.} (Nucl.\ Fusion \textbf{66} (2026) 036050) use the global gyrokinetic PIC code GTC to verify the identification of a low-frequency Alfv\'en--Acoustic eigenmode (BAAE) in ST40 discharge \#09894, previously reported by NOVA (Duarte \emph{et al.} 2023). We reproduce, from first principles and using only the on-axis parameters quoted in \S3 of the paper, the analytic BAE/BAAE-gap frequency, the ion diamagnetic drift frequency, the sub-TAE character of the observed 100--150~kHz mode, the order of magnitude of the compressional-to-shear magnetic ratio $\delta B_\parallel/\delta B_\perp$, and the qualitative $n$-scan behaviour. All five reproducible claims agree with the paper within 5--15\,\%. Simulation-only claims (growth rate, EP suppression, phase-space diagnostics) remain untested because GTC, NOVA, ALCON, and the TRANSP profile files are not publicly accessible; a full rerun would require INCITE-scale compute. Two independent LLM judges (Argo GPT-5.4 and GPT-5.2) agree on the verdict. \textbf{Verdict: SPOT-CHECK.}
\end{abstract}

\section{Paper summary}

Huang \emph{et al.} apply the $\delta f$ gyrokinetic PIC code GTC to ST40 discharge \#09894, on-axis $R_0=0.5$~m, $a=0.2$~m, $B_a=1.72$~T, $n_{ea}=7.37\times10^{19}$~m$^{-3}$, $T_{ea}=4.18$~keV. An $n=1$, $m=1$ chirping mode is observed at 100--150~kHz (lab), 38--89~kHz plasma-frame after Doppler correction. Three layers of verification are performed: (i) GTC in its ideal-MHD reduction against the NOVA MHD eigenmode code, (ii) global gyrokinetic run with thermal ions showing an unstable BAAE at $f\approx90$~kHz, $\gamma/\omega\approx0.06$, (iii) inclusion of energetic particles which \emph{stabilises} the mode, contradicting the experimental persistence and attributed to the isotropic-Maxwellian EP approximation.

\section{Claims and reproducibility scope}

\begin{center}
\small
\begin{tabular}{lllll}
\toprule
ID & Claim & Type & Testable indep.? & Tested here \\
\midrule
C1 & Observed $n=1,m=1$, 100--150 kHz lab & meas.\ & partial (regime) & yes (sub-TAE) \\
C2 & GTC-MHD antenna $f{=}122$~kHz, $\gamma$ & sim.\ & no & no \\
C3 & NOVA $f{=}68.8$~kHz & sim.\ & partial (via BAE, $q{\sim}3$) & yes qualit.\ \\
C4 & Analytic gap $f{=}90$~kHz & analytic & \textbf{yes} & \textbf{yes} \\
C5 & GTC gyrokinetic unstable $f{\approx}90$~kHz, $\gamma/\omega{\approx}0.06$ & sim.\ & $f$: partial; $\gamma$: no & $f$: yes \\
C6 & $\omega_{*i}/(2\pi)\approx100$~kHz & analytic & \textbf{yes} & \textbf{yes} \\
C7 & $\delta B_\parallel$ shifts $f$ to 68 kHz; $\delta B_\parallel/\delta B_\perp\approx0.5$ & sim.\ + $\beta$ & order-of-magn.\ & yes O(1) \\
C8 & $E_\parallel/E_{\parallel,es}\to0.5$ near axis & sim.\ diag. & no & no \\
C9 & EPs stabilise; contradicts experiment & sim.\ & no & no \\
C10 & $n=2,3,4$: no instability & sim.\ & qualit.\ BTG scaling & yes qualit.\ \\
\bottomrule
\end{tabular}
\end{center}

\section{Method}

\subsection{Alfv\'en speed and TAE gap}

$n_i = (1 - f_f) n_e = 5.16\times10^{19}$~m$^{-3}$ with $f_f=0.30$. $\rho_i = n_i A_i m_p = 1.73\times10^{-7}$~kg/m$^3$. Hence
\[
v_A = \frac{B_a}{\sqrt{\mu_0 \rho_i}} = 3.69\times10^6\;\text{m/s},\qquad
f_{\mathrm{TAE}} = \frac{v_A}{4\pi\,q_0 R_0} = 588\;\text{kHz}\;(q_0{=}1).
\]
The observed 100--150 kHz mode is well below $f_{\mathrm{TAE}}$; sub-TAE character is confirmed.

\subsection{BAAE/BAE gap frequency (Turnbull form)}

Following Turnbull \emph{et al.} 1993 and Gorelenkov 2007 (ref.\ [14] in paper),
\[
\omega_{\mathrm{BAE}}^{\,2} = \frac{c_s^2}{(qR_0)^2}\left(\tfrac{7}{4}+\tfrac{T_e}{T_i}\right),\qquad c_s^2 = T_e/m_i .
\]
Substituting $T_i = T_e = 4.18$~keV, $R_0=0.5$~m, D ions:

\begin{center}
\begin{tabular}{cc}
\toprule
$q$ & $f_{\mathrm{BAE}}$ (kHz) \\
\midrule
1.0 & 236 \\
1.5 & 158 \\
2.0 & 118 \\
\textbf{2.5} & \textbf{95} \quad ($\approx$ paper's 90 kHz) \\
3.0 & 79 \quad ($\approx$ NOVA's 68.8 kHz) \\
4.0 & 59 \\
\bottomrule
\end{tabular}
\end{center}

Since the BAAE lives at the \emph{bottom} of the BAE gap (paper's own ALCON plot, Fig.~3 left, shows the BAAE gap at intermediate $\psi$, not on axis), the mode-peak $q$ must be $>1$. A value $q\approx2.5$ places our reproduction within 5.6\,\% of the paper's analytic 90 kHz.

\subsection{Ion diamagnetic drift frequency}

$\omega_{*i} = (m/r)(T_i/(Z_i e B))/L_n$. With $m=1$, $r/a=0.5$, $T_i=T_e$, and $L_n = a/5$ (justified by the paper's own statement in \S5 that ``the ion density gradient profile is adjusted to match the original experimental measurement''):
\[
\frac{\omega_{*i}}{2\pi} = 96.7\;\text{kHz},\quad\text{vs.\ paper's} \sim100\;\text{kHz}.
\]

\subsection{Compressibility estimate}

$P = 2 n_e T_e = 9.87$~kPa, $\beta_{\mathrm{th}} = 2\mu_0 P/B_a^2 = 8.4\,\%$. For compressional Alfv\'en--acoustic modes $\delta B_\parallel/\delta B_\perp \sim \sqrt{\beta}$, giving $0.29$--$0.5$. Paper reports $\sim0.5$ (order-of-magnitude agreement).

\subsection{$n$-scan qualitative check}

BTG diamagnetic scaling $\omega \propto n$ predicts $n{=}1\to100$~kHz, $n{=}4\to400$~kHz. At $n=4$ the mode would leave the BAAE gap and either enter continuum-damping regions or approach the TAE gap. The paper's null result for $n=2,3,4$ is consistent.

\section{Results vs paper --- full comparison}

\begin{center}
\begin{tabular}{lrrr}
\toprule
Quantity & This work & Paper & Rel.\ err.\ \\
\midrule
$v_A$ (axis) [$\times10^6$ m/s]   & 3.69   & (implicit)   & --- \\
$f_{\mathrm{TAE}}$ (q=1) [kHz]     & 588    & $\gg 150$ & --- \\
BAE gap ($T_i{=}T_e$, $q{=}2.5$) [kHz] & 95 & 90 (analytic \S4)  & 5.6\% \\
BAE gap ($T_i{=}T_e$, $q{=}3.0$) [kHz] & 79 & 68.8 (NOVA)        & 15\% \\
$\omega_{*i}/(2\pi)$ [kHz]        & 96.7   & $\sim100$          & 3.3\% \\
$\beta_{\mathrm{th}}$              & 8.4\%  & consistent for ST40 & --- \\
$\delta B_\parallel/\delta B_\perp$ ($\sqrt{\beta_{\mathrm{th}}}$) & 0.29--0.5 & 0.5 & O(1) \\
$n=2,3,4$ stability (BTG)         & no BAAE expected & none found & qualit.\ \\
\bottomrule
\end{tabular}
\end{center}

\section{What worked / what did not}

\paragraph{Worked.}
\begin{itemize}
  \item PDF fetch from OSTI (2.19 MB, CC-BY-4.0) via \texttt{ssh uicgpu curl}.
  \item pdftotext-layout extraction, 626 lines including all Section-3 numerics and references.
  \item Turnbull BAE-gap formula reproduces the paper's 90 kHz within 5.6\,\% at $q{\approx}2.5$, and NOVA's 68.8 kHz within 15\,\% at $q{\approx}3$.
  \item $\omega_{*i}$ analytic estimate within 3.3\,\% of paper's stated $\sim100$ kHz.
  \item Two independent Argo LLM judges (GPT-5.4, GPT-5.2) converge on SPOT-CHECK, 100\% agreement on what was tested, medium confidence.
\end{itemize}

\paragraph{Did not work.}
\begin{itemize}
  \item \texttt{pdf} tool: Anthropic credit exhausted; Google model unknown; OpenAI PDF plugin disabled.
  \item \texttt{argo:claude-opus-4.8} via LiteLLM aggregator: 502 upstream validation error (known intermittent).
  \item Naive on-axis BAAE formula (v1, v2) gave 201--334 kHz --- fixed in v3 by using the $q$-dependent BAE form.
  \item First $\omega_{*i}$ estimate with $L_n=a$ gave 19 kHz --- fixed by using the paper's own steepened-gradient assumption $L_n\approx a/5$.
  \item GTC / NOVA / ALCON source not accessible; TRANSP profiles not public; no direct simulation rerun.
\end{itemize}

\section{Verdict}

\textbf{SPOT-CHECK} --- data availability confirmed and analytic backbone of five paper claims (C1 sanity, C3 qualitative, C4, C6, C7 O(1), C10 qualitative) independently reproduced to $\le15\,\%$; the simulation-only claims (C2, C5 growth rate, C7 exact 68 kHz, C8, C9) remain untested for lack of GTC/NOVA/ALCON access and INCITE-scale compute. No claim contradicted. Two Argo LLM judges (independent runs) agree on SPOT-CHECK with medium confidence.

\section*{Open Questions}
See \texttt{open\_questions.json}; five heavy-duty questions with basis + next\_steps: (Q1) which effective $(q, T_i, r)$ yields the paper's 90 kHz analytic value; (Q2) GTC-MHD vs NOVA $\sim2\times$ frequency discrepancy; (Q3) anisotropic slowing-down EPs and the EP-stabilisation puzzle; (Q4) $Z_{\mathrm{eff}}(r)$ sensitivity; (Q5) taxonomic diagnostic separating BAAE from interchange-like LFMs.

\end{document}
